\section{The conditioned edge measure and the Gaussian inputs}
\label{sec:edge-law}

The limiting state space is determined by the upper spectral edge of the
interaction matrix.  The spherical model makes this state space explicit:
after diagonalization, its only remaining interaction is a constraint on the
sum of the squared coordinates.  We construct the resulting probability measure on infinitely many coordinates, following the equilibrium
edge-conditioning description of Du--Huang \cite[Section~2]{DuHuang2026Overlap}.  We then explain the change of disorder law that transfers spherical
calculations to the cube, and record the cap estimates used in the dynamical
argument.

\subsection{The heat-bath form and the edge coordinates}

Write $\mu=\pi_{n,J}$ and put $T=n^{2/3}$.  For a function on the cube, let
\[
 D_i f(X)=\frac{f(X^{i,+})-f(X^{i,-})}{2},\qquad
 t_i(X)=\tanh h_i(X),\qquad w_i(X)=1-t_i(X)^2,
\]
where $X^{i,\pm}$ is obtained by setting the $i$th spin to $\pm1$ and
$h_i(X)=n^{-1/2}\sum_{j\ne i}J_{ij}X_j$.  The nonnegative heat-bath
operator and its Dirichlet form are
\begin{equation}\label{eq:edge-heatbath-form}
 H=\sum_{i=1}^n(X_i-t_i)D_i,
 \qquad
 \cE(f,g)=\mu\!\left[\sum_{i=1}^n w_iD_ifD_ig\right],
 \qquad \cE(f)=\cE(f,f).
\end{equation}
Thus $P_t^J=e^{-tH}$.  Forms and inner products of vector-valued functions
include summation over the vector index.  In particular, the clock
normalization in \eqref{eq:edge-heatbath-form} is rate one at each site.

Until Section~\ref{sec:comparison}, the couplings are Gaussian.  Let $W$ have
off-diagonal entries $J_{ij}/\sqrt n$, and add independent diagonal entries
with law $N(0,2/n)$.  The extra energy is constant on the cube, so it changes
neither $\mu$ nor $H$.  The completed matrix is GOE.  Diagonalize it as
\[
 W=U\Lambda U^{\mathsf T},\qquad
 \lambda_1\ge\cdots\ge\lambda_n,
 \qquad d_a=2-\lambda_a,\qquad g_{a,n}=n^{2/3}d_a.
\]
After independent choices of eigenvector signs, $U$ is Haar orthogonal
conditionally on $\Lambda$.  Its columns are denoted by $v_a$.  Set
\[
 Y_a=v_a^{\mathsf T}X,\qquad x_a=n^{-1/3}Y_a,
 \qquad
 B_n(z)=\sum_{a=1}^n\langle x_a,(z+TH)^{-1}x_a\rangle_\mu,
 \quad z>0.
\]
The trace $B_n(z)$ is independent of the chosen orthonormal basis and is
the Laplace transform of the rescaled autocorrelation:
\[
 B_n(z)=\int_0^\infty e^{-zs}A_{n,J}(s)\,ds.
\]
We use $x_a$ also for the coordinate functions of the limiting measure.

Let $\rho_{\rm sc}$ be the semicircle probability measure
\[
 d\rho_{\rm sc}(\lambda)=\frac1{2\pi}
 \sqrt{4-\lambda^2}\,\1_{[-2,2]}(\lambda)\,d\lambda.
\]
The GOE edge convergence theorem gives the joint convergence of each fixed
number of $g_{a,n}$ to a simple, increasing point process
$g=(g_1,g_2,\ldots)$; see \cite{RamirezRiderVirag2011}.  This definition fixes the sign and
scale of the edge process throughout the paper.

\begin{lemma}[Spectral bounds]\label{lem:edge-spectral}
The GOE eigenvalues have the following additional properties.
\begin{enumerate}
\item $\max_a|\lambda_a|\le3$ with probability tending to one.
\item There are tight random bounds $K_n\ge0$ and $C_n>0$ such that
$g_{a,n}\ge C_n^{-1}a^{2/3}$ for every $a>K_n$.
\item For every $\epsilon>0$,
\begin{equation}\label{eq:edge-trace-double-limit}
 \lim_{L\to\infty}\limsup_{n\to\infty}
 \mathbb P\!\left(
 g_{1,n}+L\le0\ \text{or}\
 \left|n^{1/3}-\sum_{a=1}^n\frac1{g_{a,n}+L}-\sqrt L\right|
       >\epsilon\right)=0.
\end{equation}
\end{enumerate}
In the event appearing in the last formula, the sum is evaluated only when
$g_{1,n}+L>0$.
\end{lemma}

\begin{proof}
We use the eigenvalue-only spectral estimate of \cite[Proposition~3.1]{CS}.  For every
deterministic sequence $L_n\to\infty$ with $L_n\le\log n$, that proposition
controls both resolvent traces for
$s\in[L_nn^{-2/3},5]$, with positive denominators.  At
$s=L_nn^{-2/3}$, multiplying its first-trace error by $n^{1/3}$ gives
\[
 O\!\left(L_n^{-1}+L_n^{-1/8}
       +L_n^{-1/4}\sqrt{\log(2+L_n)}\right)=o(1).
\]
The semicircle resolvent equals
$(2+s-\sqrt{s(4+s)})/2$, whose expansion supplies the main term
$\sqrt{L_n}$ and an additional error $O(L_nn^{-1/3})=o(1)$.

The second-trace estimate yields, for the same interval of $u$,
\[
 \#\{a:d_a\le u\}
 \le 4u^2\sum_a(d_a+u)^{-2}
 \le Cnu^{3/2}.
\]
Indeed $0<d_a+u\le2u$ for every counted eigenvalue.  It follows that
$g_{a,n}\ge c a^{2/3}$ once $a>CL_n^{3/2}$.  The uniform norm bound is
part of the same proposition.

These conclusions hold for every arbitrarily slowly divergent $L_n$.
If \eqref{eq:edge-trace-double-limit} failed, one could select
$L_j\to\infty$ and $n_j\to\infty$, with $\log n_j\ge L_j$, along which
the failure probability stayed positive, and embed the values
$L_{n_j}=L_j$ in a sufficiently slow divergent sequence.  This contradicts
the preceding estimate.  The identical criterion, applied to the least
index after which the counting consequence holds, makes that index tight.
This proves the second assertion as well.
\end{proof}

\begin{remark}[Convergence on sets of large disorder probability]
\label{rem:edge-tight-convergence}
We will use constants that are bounded on disorder events of probability
arbitrarily close to one.  A bound of this kind is called a tight bound.
For subsequential arguments, include these bounds, the eigenvalues, and the
countably many relevant error quantities in the joint law before taking a
Skorokhod representation.  On the represented space, finite bounds and
coordinate convergence can then be used deterministically.  The conditional
laws of the remaining finite-dimensional variables lift this representation
to the original finite problems.  A further diagonal extraction handles
quantities that vanish in an iterated limit such as
\eqref{eq:edge-trace-double-limit}.  The representation itself makes no assertion of convergence for the
original sequence: one must identify every subsequential limit.  Once this
identification is made, the subsequence criterion gives the asserted joint
convergence in distribution.  In particular, if every subsequence of an
intrinsic error quantity has a further represented subsequence tending to
zero, that quantity tends to zero in probability on the original disorder
space.  No uniform integrability on exceptional disorders is required.

A separate uniformity convention is essential for variational arguments.
When an estimate is applied to a disorder-dependent test, its quenched
inequality is understood to hold on one event for every test in the stated
class, with the displayed bounds on its norm, energy, and polynomial growth.
One way to obtain such an event, used below, is to bound the error for every
test by its stated norms times a common nonnegative random error $R_n$.
Convergence $R_n\to0$ then permits arbitrary measurable choices of the test.
Convergence on a countable core alone is not used to infer this uniformity.
\end{remark}

\subsection{Conditioning the infinite Gaussian family}

Choose an admissible parameter $b$ with $b+g_1>0$, and write
$l_a=(g_a+b)^{-1}$.  A convenient measurable choice is
$b=1+(-g_1)_+$; use the corresponding $b_n=1+(-g_{1,n})_+$ in finite
dimension.  By Lemma~\ref{lem:edge-spectral},
\[
 l_{a,n}=(g_{a,n}+b_n)^{-1}\longrightarrow l_a,
 \qquad
 \lim_{m\to\infty}\sup_n\sum_{a>m}l_{a,n}^2=0
\]
on the represented sets just described, after discarding finitely many
initial $n$.  In particular $\sum_a l_a^2<\infty$.

For independent $Z_a\sim N(0,l_a)$, the centered series
$S=\sum_a(Z_a^2-l_a)$ converges almost surely and in $L^2$.  Define
\begin{equation}\label{eq:source-1}
 D(b)=\lim_{L\to\infty}\left\{
 \sqrt L-\sum_{a\ge1}\left(\frac1{g_a+b}-\frac1{g_a+L}\right)\right\}.
\end{equation}
Here $L$ is restricted to values with $L+g_1>0$.  The limit is in probability
with respect to $g$; a deterministic subsequence gives an almost sure
version.  The next proposition both justifies the limit and gives a
pointwise meaning to conditioning on $S=D(b)$.

\begin{proposition}[The conditioned edge measure]
\label{prop:edge-conditioned-law}
Formula~\eqref{eq:source-1} defines a finite random variable, and
\[
 D_n(b_n):=n^{1/3}-\sum_{a=1}^n l_{a,n}
       \ \Longrightarrow\ D(b)
\]
jointly with the edge coordinates.  Almost surely $\sum_a l_a=\infty$.
The law of $(Z_a)_{a\ge1}$ conditioned on $S=D(b)$ has a version
$\Pi=\Pi_g$ characterized by continuous conditional densities.  It is
independent of the admissible choice of $b$.

Let $\mu^{\rm sp}$ denote normalized surface measure on the sphere of
radius $\sqrt n$, tilted by $\exp(X^{\mathsf T}WX/2)$.  The fixed-coordinate marginals of $n^{-1/3}U^{\mathsf T}X$ under
$\mu^{\rm sp}$ converge jointly in distribution with the eigenvalues to
the random marginals of $\Pi$.  On the spectral representations of
Remark~\ref{rem:edge-tight-convergence}, this is ordinary weak convergence
of measures, including convergence of every fixed polynomial moment.  On the sets in Remark~\ref{rem:edge-tight-convergence},
for every fixed $0\le j<\infty$ one has
\begin{equation}\label{eq:edge-coordinate-moments}
 \int |x_a|^j\,d\Pi\le C_j l_a^{j/2},
 \qquad
 \mu^{\rm sp}|x_a|^j\le C_j l_{a,n}^{j/2},
\end{equation}
uniformly in the retained coordinate $a$.

For any sufficiently large fixed $m$, the tail $(x_a)_{a>m}$ under
$\Pi$ has a density relative to $\bigotimes_{a>m}N(0,l_a)$.  Conditional
on this tail, the first $m$ coordinates lie on the sphere of squared radius
\begin{equation}\label{eq:edge-polar-radius}
 D(b)+\sum_{a\le m}l_a-\sum_{a>m}(x_a^2-l_a),
\end{equation}
with the Gaussian angular density on that sphere.  In particular,
$\sum_a(x_a^2-l_a)=D(b)$ holds $\Pi$-almost surely.
\end{proposition}

\begin{proof}
For fixed admissible $L$, the identity
\[
 D_n(b_n)=n^{1/3}-\sum_a\frac1{g_{a,n}+L}
 -\sum_a\left(\frac1{g_{a,n}+b_n}-\frac1{g_{a,n}+L}\right)
\]
separates the divergent trace from an absolutely convergent difference.
The tail of the second sum is bounded by a constant times
$|L-b_n|\sum_{a>m}a^{-4/3}$ on each set of bounded spectral data, so that
sum passes to the limit.  Lemma~\ref{lem:edge-spectral} then shows that the
expressions in braces in \eqref{eq:source-1} are Cauchy in probability,
and that $D_n(b_n)$ has the asserted joint limit.  If $\sum_a l_a$ were
finite, those expressions would tend to $+\infty$ as $L\to\infty$;
this contradicts the finite limit.  Changing $b$ changes $D(b)$ by the
corresponding absolutely convergent resolvent difference.

We give the density argument explicitly.  Let $q$ be the density of $S$,
and let $q_m$ be the density of $\sum_{a>m}(Z_a^2-l_a)$.
After any fixed coordinates are omitted, a sufficiently large fixed group
of remaining coordinates supplies an integrable bound for the characteristic
function: the modulus of the factor for coordinate $a$ is
$(1+4t^2l_a^2)^{-1/4}$.  Fourier inversion therefore gives continuous,
bounded densities.  The same bound applies to the finite sums and proves
uniform convergence of their densities when the variances converge and
their squared tails vanish.

The centered sum after any finite deletion has support all of $\R$.
Its unbounded right support comes from one square.  To reach an arbitrarily
negative interval, make a sufficiently long finite block of squares small,
using $\sum l_a=\infty$, and keep the remaining centered sum in a bounded
interval of positive probability.  Convolution with a finite block whose
square-sum density is positive in the interior of its support now gives
$q(s)>0$ and $q_m(s)>0$ for every $s\in\R$.  Thus the first $m$
coordinates of the required conditional law have density
\begin{equation}\label{eq:edge-cylinder-density}
 \frac{q_m\!\left(D(b)-\sum_{a\le m}(x_a^2-l_a)\right)}{q(D(b))}
 \prod_{a\le m}\frac{e^{-x_a^2/(2l_a)}}{\sqrt{2\pi l_a}}.
\end{equation}
Convolution shows that these densities are consistent.  They depend
measurably on the spectral data, so their extension is a random probability
measure $\Pi_g$.

For completeness, they admit a disintegration on the constraint itself.
Let $q^{\rm head}_m$ be the density of
$\sum_{a\le m}(Z_a^2-l_a)$.  Give the tail product Gaussian law the
weight
\[
 \frac{q^{\rm head}_m\!\left(D(b)-\sum_{a>m}(x_a^2-l_a)\right)}{q(D(b))}.
\]
Given the tail, use polar disintegration of the first $m$ Gaussian
coordinates at the radius \eqref{eq:edge-polar-radius}.  This produces
\eqref{eq:edge-cylinder-density}, has an absolutely continuous tail law,
and enforces the constraint almost surely.

In finite dimension, conditioning independent normals of variances
$l_{a,n}$ on $\sum_a x_a^2=n^{1/3}$ gives exactly the spherical Gibbs
measure in the stated coordinates: the term involving $b_n$ in the Gaussian
exponent is constant on that sphere.  The density convergence above and
$q(D(b))>0$ imply convergence of each cylinder marginal.  Their densities
are bounded by a constant times the corresponding product Gaussian
density.  To retain an arbitrary single coordinate uniformly, choose the
integrable Fourier bound from a fixed group of leading coordinates other
than that coordinate.  This proves \eqref{eq:edge-coordinate-moments}
and uniform integrability of all fixed polynomial moments.

The argument applies to every admissible $b$ on sets where its distance
from the pole is bounded below and its magnitude is bounded above; these
sets exhaust the admissible choices.  The finite conditioned spherical law
is independent of $b_n$, so its limiting cylinder marginals are also
independent of $b$.  The conclusion holds simultaneously for all
admissible $b$.  First take a common probability-one event for rational
parameters.  On every compact admissible interval, the resolvent-difference
series converges uniformly, so $D(b)$ is continuous.  The squared variance
tails and the Fourier majorants above are also uniform on that interval.
Thus the conditional cylinder densities in
\eqref{eq:edge-cylinder-density}, and their integrals against bounded
continuous cylinder tests, depend continuously on $b$.  Equality for
rational parameters therefore extends to the whole interval.
\end{proof}

\begin{lemma}[The bulk spectral weight]\label{lem:edge-bulk-weight}
For a spherical Gibbs spin, the random probability measure
\[
 \frac1n\sum_a Y_a^2\delta_{\lambda_a}
\]
converges in probability, conditionally on the spectral data as above, to
$(2-\lambda)^{-1}d\rho_{\rm sc}(\lambda)$.
\end{lemma}

\begin{proof}
First test against a bounded continuous function supported away from $2$.
Before conditioning, the Gaussian variance of its integral is $O(n^{-1})$,
and its mean converges to the stated integral by the semicircle law.
Integrate out a fixed number of leading Gaussian coordinates when imposing
the norm constraint.  Their square-sum density gives a bounded density
ratio, so the concentration persists after conditioning.  Each of those
leading coordinates has mass $n^{-1/3}x_a^2=o(1)$ in probability by
\eqref{eq:edge-coordinate-moments}.  Finally both the finite measure and
its proposed limit have mass one, since
$\int(2-\lambda)^{-1}d\rho_{\rm sc}(\lambda)=1$.  This extends convergence
across the edge.  The same argument permits bounded tests with cuts at
continuity points in the bulk.
\end{proof}

\subsection{The limiting autocorrelation and its time constant}

For a smooth compactly supported cylinder function $F$ on $\R^{\N}$,
consider the quadratic form
\[
 \cE_\Pi(F)=\int\sum_a|\partial_aF|^2\,d\Pi.
\]
Proposition~\ref{prop:gradient-domain} will prove that this form is
closable and identify its domain.  Denote the nonnegative operator associated
to its closure by $H_\Pi$.  The law in Theorem~\ref{thm:main} is specified by
\begin{equation}\label{eq:source-2}
 \cA(s)=\sum_{a\ge1}
 \langle x_a,e^{-s c_*H_\Pi}x_a\rangle_\Pi,
 \qquad s>0,
\end{equation}
where $c_*>0$ is deterministic.  Finiteness, continuity, and the remaining
assertions about this random function are proved in
Sections~\ref{sec:trace} and~\ref{sec:functional}.

The coefficient is $c_*=\omega(\{0\})$, where $\omega$ is the
deterministic finite positive measure of microscopic heat-bath energies
characterized by the moment prescription \eqref{eq:source-3} in
Section~\ref{sec:microscopic}.  Proposition~\ref{prop:micro-moment-measure}
constructs this measure and proves its determinacy;
Theorem~\ref{thm:micro-relaxation} proves that its atom at zero is strictly
positive and is the coefficient in \eqref{eq:source-2}.

\subsection{Haar bias and Gaussian cap observations}

We now connect the spherical calculations with the cube.  All partition
functions use probability reference measures: uniform probability on the
cube and normalized surface measure on the sphere.  Conditional on the
spectrum, put
\[
 L_\circ(U)=\frac{Z_{\rm cube}(W)}{Z_{\rm sp}(W)}.
\]
Haar averaging gives $E_U L_\circ=1$.  Define the probability law
$\mathbf P$, and its expectation $\mathbf E$, by drawing $U$ with
Haar density $L_\circ$, then drawing a cube spin with law $\mu$.
Any observation noises introduced below are independent of this spin
conditionally on the matrix.  The spectrum is held fixed in
$\mathbf P$ and $\mathbf E$.

\begin{lemma}[Spherical identity under Haar bias]
\label{lem:edge-biased-law}
Under $\mathbf P$, $U^{\mathsf T}X$ has exactly the spherical Gibbs law
in spectral coordinates.  After the row sign change sending $X$ to
$\mathbf1$, conditional on $y=U^{\mathsf T}X$ the eigenframe is Haar
subject to $Uy=\mathbf1$.  Moreover, on the eigenvalue events used below,
\begin{equation}\label{eq:edge-uncapped-comparison}
 E_U(L_\circ-1)^2=o(1).
\end{equation}
\end{lemma}

\begin{proof}
In the joint weighted law, $Z_{\rm cube}$ cancels the Gibbs denominator.
For each fixed cube vertex, Haar rotation sends that vertex to uniform
surface measure on the sphere of radius $\sqrt n$.  The remaining energy
weight is $\exp(y^{\mathsf T}\Lambda y/2)$, proving the first assertion.
The same invariance gives the conditional Haar law on the stated coset.
Finally, the uncapped zero-field calculation in the proof of the static
cap comparison of \cite{CM}, Section~15, paragraph ``Zero field and
overlaps bounded away from zero,'' gives
$E_UZ_{\rm cube}^2/(Z_{\rm sp})^2=1+o(1)$.
Together with $E_UL_\circ=1$, this is
\eqref{eq:edge-uncapped-comparison}.
\end{proof}

For example, if $G\ge0$ is a spin or observation functional, then
\[
 E_U\!\left[\1_{\{L_\circ\ge1/2\}}Q_\mu G\right]
       \le2\mathbf E G.
\]
Here $Q_\mu$ denotes the ordinary joint law of the Gibbs spin and its
observations at fixed disorder.  This inequality and Markov's inequality
transfer biased mean estimates to quenched estimates; the discarded disorder
event has probability $o(1)$ by \eqref{eq:edge-uncapped-comparison}.

The cap at scale $p>0$ is defined by
\[
 P_p=\1_{\{d\le p^2\}},\quad \cH_p=\operatorname{ran}P_p,\quad
 B_p=(p^2-d)_+,\quad W_p=W-B_p,\quad k_p=np^3,\quad M_p=\sqrt{np}.
\]
Matrix functions in these formulas are conjugated by $U$ when used in
physical coordinates.  Figure~\ref{fig:spectral-cap} shows how the cap
flattens the upper edge and how the removed part becomes observation precision.
\begin{figure}[tb]
\centering
\begin{tikzpicture}[x=1.65cm,y=1.2cm,>=Stealth,font=\small]
 \fill[blue!8] (2,2)--(3.7,3.7)--(3.7,2)--cycle;
 \draw[->] (0,0)--(4.25,0) node[right] {$\lambda$};
 \draw[->] (0,0)--(0,4.05) node[above] {eigenvalue};
 \draw[gray,dashed,thick] (0.35,0.35)--(3.85,3.85);
 \draw[blue!60!black,very thick] (0.35,0.35)--(2,2)--(3.9,2);
 \draw[dotted] (2,0)--(2,2);
 \draw[dotted] (0,2)--(2,2);
 \node[below] at (2,0) {$2-p^2$};
 \node[left] at (0,2) {$2-p^2$};
 \draw (3,0.045)--(3,-0.045) node[below] {$2$};
 \node[gray,rotate=36,above] at (1.1,1.1) {original eigenvalue};
 \node[blue!60!black,below] at (3.15,1.9) {$\min(\lambda,2-p^2)$};
 \draw[<->,thick] (3.5,2.06)--(3.5,3.45);
 \node[right,align=left] at (3.57,2.8)
   {observation\\precision};
 \draw[<->] (2,-0.65)--(3.95,-0.65);
 \node[below,align=center] at (2.98,-0.65)
   {active eigendirections\\$\lambda>2-p^2$};
\end{tikzpicture}
\caption{A spectral cap, shown near the upper edge. The posterior interaction
$W_p$ retains eigenvalues below $2-p^2$ and flattens those above it.
The removed amount $(\lambda-(2-p^2))_+$ is the precision of the Gaussian
observation in that eigendirection. The drawing is schematic; its geometry
illustrates the defining spectral functions, not an eigenvalue density.}
\label{fig:spectral-cap}
\end{figure}


Observe
$h_p=B_pX+\mathsf N(0,B_p)$, with independent Gaussian increments as
$B_p$ increases, starting from zero precision if necessary.  Bayes' formula
makes the posterior interaction $W_p$ and the posterior field $h_p$.
Write $\mu_{p,h}$ for the cube posterior at a specified field $h$.  Let $\phi_p(h)$ be its log partition function,
$m_p(h)=\nabla\phi_p(h)$ its mean, and
$K_p(h)=P_p\nabla^2\phi_p(h)P_p$ its covariance restricted to $\cH_p$.
Superscript ${\rm sp}$ denotes the corresponding spherical quantities.
At path fields we often suppress the argument $h_p$.

On the joint space of the eigenframe and the spectral observation path,
the density of $\mathbf P$ relative to independent Haar measure and the
spherical observation law, through scale $p$, is
\begin{equation}\label{eq:source-4}
 L_p=\exp\{\phi_p(h_p)-\phi_p^{\rm sp}(h_p)\}.
\end{equation}
Indeed the likelihood of the observations is a common Gaussian factor
multiplied by the posterior partition function divided by the original one;
the factor $L_\circ$ cancels the two original partition functions.
The spin in spectral coordinates already has the spherical law by
Lemma~\ref{lem:edge-biased-law}; its observation noises are independent
Gaussian increments with covariance determined by the spectrum.  Hence the
entire spectral observation path under $\mathbf P$ has exactly its
spherical law.  We normalize its coordinates by
\[
 u_{p,a}=\frac{(U^{\mathsf T}h_p)_a}{p^2n^{1/3}},
 \qquad r=n^{-1/3+\delta},\qquad \delta=10^{-4}.
\]
We will use $r\le p\le p_{\rm top}$, allowing fixed-factor buffers at the
endpoints.  The values $d_a=p^2$ cause no discontinuity in the observations;
they may be omitted in scale integrals, and zero-variance coordinates are
omitted from density formulas.

\subsection{Saddles and cap estimates}

For a field $h\in\cH_p$, its spherical saddle $\alpha=\alpha_p(h)$ is
the unique solution above the pole of
\begin{equation}\label{eq:source-5}
 R_{p,\alpha}=(\max(d,p^2)+\alpha)^{-1},\qquad
 n=\Tr R_{p,\alpha}+\|R_{p,\alpha}h\|^2.
\end{equation}
The annulus $|\alpha_p(h)|\le s_0p^2$ is separated from zero field.  Its
radius as a function of $\alpha$ is
\[
 \ell_p(\alpha)=(p^2+\alpha)\sqrt{n-\Tr R_{p,\alpha}}.
\]
The following statement collects the precise Gaussian inputs needed later.
All upper cutoffs and tolerances in it are fixed constants before $n$ tends
to infinity.

\begin{proposition}[Cap inputs]\label{prop:edge-cap-inputs}
Fix a finite field radius $R$ and positive fixed error tolerances.  There
are sufficiently small constants $p_{\rm top},s_0>0$ such that the following
hold with probability tending to one in the range
$np^3\ge c\sqrt{\log n}$ and $p\le p_{\rm top}$.

\emph{Spectral and spherical estimates, at every scale.}
The spectral dimension satisfies $\dim\cH_p\asymp k_p$,
$d_1\ge-o(p^2)$ uniformly in the cap range, and
$\Tr R_{p,0}^2=O(n/p)$.  On $|\alpha|\le s_0p^2$,
$\ell_p(\alpha)\asymp p^2M_p$ and
$\ell_p'(\alpha)\asymp M_p$.  For
$\|h\|\le Rp^2M_p$, the saddle precision has gap comparable to $p^2$.
At $\alpha=\alpha_p(h)$, the density at $n$ of the squared norm of
$\mathsf N(R_{p,\alpha}h,R_{p,\alpha})$ is comparable to $\sqrt{p/n}$.
After multiplication by $e^{v^{\mathsf T}X}$ and renormalization, the
Gaussian mean becomes $R_{p,\alpha}(h+v)$ while its covariance remains
$R_{p,\alpha}$; its squared-norm density at $n$ is at most a fixed
multiple of $\sqrt{p/n}$.  For $m$ replicas, where $m$ is fixed, consider
rowwise linear or quadratic tilts that preserve Gaussianity.  The joint
density of their squared norms is at most $C_m(p/n)^{m/2}$ under the
following hypotheses.  The spectral rows remain independent after the rowwise tilt,
the tilted precision is positive definite, and on at least $ck_p$ rows
the post-tilt row covariance $C_a$ satisfies
$c p^{-2}\Id_m\preceq C_a\preceq C p^{-2}\Id_m$.
All comparison constants here are fixed.

\emph{Curvature estimates, at every scale.}
On the annulus, for any prescribed fixed $\epsilon_0>0$,
\begin{equation}\label{eq:source-6}
 K_p\preceq(1+\epsilon_0)p^{-2}\Id_{\cH_p},\qquad
 p^2(K_p)_{ee}\le1-\epsilon_1,
 \quad e=h/\|h\|,\quad \epsilon_1=0.1.
\end{equation}
On the entire field ball one also has $K_p\preceq C_Rp^{-2}\Id_{\cH_p}$,
where $C_R$ may be chosen before $\epsilon_0$.

\emph{Static comparisons and cube projection tails, on geometric lists.}
For any prescribed fixed geometric list of scales, uniformly over the field ball,
\[
 |\phi_p(h)-\phi_p^{\rm sp}(h)|\le\epsilon_v k_p,
 \qquad
 \|P_pm_p(h)-R_{p,\alpha_p(h)}h\|\le\epsilon_gM_p.
\]
Any fixed finite family of such lists is allowed.  For each
$h\in\cH_p$ with $\|h\|\le Rp^2M_p$, the cube posterior $\mu_{p,h}$
satisfies, uniformly for every $v\ge v_0$ with a sufficiently large fixed
threshold $v_0$,
\[
 \mu_{p,h}\{\|P_pX\|>vM_p\}\le C\exp(-ak_pv^6).
\]
For $v>p^{-1/2}$ the event is empty, since $\|X\|=\sqrt n$.
Here $a>0$, and the same bound holds under the uncapped zero-field cube
measure $\mu$ at the same list scales.

The spherical projection-tail estimate holds at every scale in the stated
range, for capped bounded-field and uncapped zero-field measures.  Also, at
every such cap scale, for independent
spherical replicas at a common field in the prescribed ball,
\[
 (\mu_{p,h}^{\rm sp})^{\otimes2}
 \{|n^{-1}X^{\mathsf T}X'|>\rho\}
 \le C\exp(-cn\rho^3),\qquad Cp\le\rho\le\rho_0,
\]
for a sufficiently small fixed $\rho_0>0$.  A fixed finite number of
replicas obeys the corresponding bound for any pair.
\end{proposition}

\begin{proof}[Source statements and the radial consequence]
The spectral, saddle, and density assertions are the conclusions of
\cite{CM}, Sections~15.2--15.3, ``Spectral estimates at the bottom
scale'' and ``Spherical saddles, norm densities, and tails.''  The static
comparison and its projection tails are Proposition~15.1 in that section.  Its order of choices is the one stated here:
fix the radius, errors, and geometric ratios, then decrease the upper cutoff.
The overlap estimate is the spherical estimate preceding that proposition.

The upper covariance bounds are Theorem~18.6, ``Uniform curvature of the true
cap potentials,'' in \cite{CM}.  Its lower endpoint is
$(\sqrt{\log n}/n)^{1/3}$, with fixed-factor changes permitted, so it
covers our scale $r$ and its fixed buffers.  We also use a consequence
retained in that theorem's proof.  Let $\gamma$ be the fixed geometric ratio chosen in that proof, and set
$R_{p,\alpha}^H=(P_pR_{p,\alpha}P_p)|_{\cH_p}$.  At an intermediate
evaluation ratio $x\in[\gamma^2,\gamma]$, Lemma~16.2, ``Evaluation with the induction
extension,'' in \cite{CM} bounds the radial entry of
$p^2(K_p-R_{p,\alpha}^H)$ above by $-0.26$.  The theorem's proof chooses
$s_0$ so that every target saddle lies inside the parent agreement annulus
with fixed slack, and applies exactly this deterministic evaluation.
Therefore
\[
 p^2(K_p)_{ee}\le\frac1{1+\alpha/p^2}-0.26
                  \le\frac1{1-s_0}-0.26\le0.9
\]
after decreasing $s_0$.  This proves the second part of
\eqref{eq:source-6} with a fixed tolerance.
\end{proof}

The list restriction in the static comparison is harmless, but its removal
requires a short argument because later observation identities integrate
over all scales.

\begin{lemma}[Interpolation between cap scales]
\label{lem:edge-static-interpolation}
The value, projected-mean, and cube projection-tail conclusions of
Proposition~\ref{prop:edge-cap-inputs} hold at every scale in the stated
range, after prescribing slightly smaller errors on a sufficiently fine
fixed geometric list and decreasing the upper cutoff.
\end{lemma}

\begin{proof}
First use an ordinary fixed-ratio list.  Monotonicity of $P_p$ and a
fixed-factor change in $M_p$ extend the uncapped projection tail to every
$p$.  Removing $B_p$ from the interaction changes its normalizing factor
by
\[
 \mu\exp(-X^{\mathsf T}B_pX/2)\ge
 \mu\{\|P_pX\|\le VM_p\}\,
       \exp(-\|B_p\|V^2M_p^2/2)\ge e^{-Ck_p},
\]
where $V$ is fixed and large.  We used $\|B_p\|=O(p^2)$ on the spectral
event.  Thus the projection tail persists at cap zero field.  A field
$\|h\|\le Rp^2M_p$ adds at most $Rk_pv$ on the shell
$\|P_pX\|=vM_p$; this is dominated by $ak_pv^6$ for large $v$.
The zero-field measure is even, so its linear-tilt normalizer is at least
one.  This gives the same tail on the field ball, with an adjusted
starting threshold.  The argument also bounds $\|P_qX\|/M_q$ under the
$p$-cap posterior whenever $q\ge p$ and $q/p$ is bounded: use the uncapped
$P_q$ tail, the same denominator lower bound, and $P_ph=h\in\cH_q$.
This is the tail needed to compare two nearby caps.  The spherical estimates give the identical conclusion.

Now choose a fixed geometric list with consecutive ratios close to one,
and let $q$ be its first scale above $p$.  Then $\cH_p\subset\cH_q$ and
\[
 \|W_p-W_q\|\le C(1-p/q)q^2,
 \qquad \operatorname{ran}(W_p-W_q)\subset\cH_q.
\]
For $h\in\cH_p$ in the prescribed ball, restrict first to
$\|P_qX\|\le VM_q$.  The change in log weight there is at most
$C(1-p/q)V^2k_q$.  On its complement, the projection tail controls the
exponential of the quadratic change: the exponent
$C(1-p/q)k_qv^2$ is dominated by the negative $ak_qv^6$ tail.
Choose $V$ first and the list ratio second.  The resulting log-partition
change is as small as desired in units $k_q$, for both cube and sphere.
Their value comparison at $q$ therefore implies the value comparison at
$p$, since $k_p/k_q$ is arbitrarily close to one.

For the mean, use the value comparison on a slightly larger field ball.
The frozen Gaussian saddle identity and its norm-density bound give, for
any unit $v\in\cH_p$ and both signs,
\[
 \log\mu_{p,h}e^{\pm t v^{\mathsf T}(X-R_{p,\alpha}h)}
       \le 2\epsilon_v k_p+\tfrac12t^2v^{\mathsf T}R_{p,\alpha}v+C.
\]
Jensen's inequality, with $t=\theta p^2M_p$, bounds the projected mean
discrepancy divided by $M_p$ by
$C\theta+(2\epsilon_v+C/k_p)/\theta$.  Choose $\theta$ from the required
mean tolerance and then $\epsilon_v$ sufficiently small.  This proves the
claim uniformly at every intermediate scale.
\end{proof}

The estimates of this section supply the ordinary Gaussian caps and the
spherical description of their observations.  The next section derives the
quantitative comparison and score estimates needed to retain the exact
$n^{2/3}$ time scale.
